\section{Kết luận}

% Slide 21
\begin{frame}[t]{Kết luận}
  \begin{itemize}
    \setlength{\itemsep}{8pt}
    \item \textbf{Dữ liệu và phân tích:}
    \begin{itemize}
      \item Chuẩn hóa thành công bộ dữ liệu PAPI từ 14 tệp Excel thô (2011\textendash 2024).
      \item Triển khai 5 hướng phân tích đa chiều: tổng quan, thời gian, không gian, lĩnh vực và phân nhóm.
    \end{itemize}
    \item \textbf{Trực quan hóa và hệ thống:}
    \begin{itemize}
      \item Xây dựng Dashboard bằng React và FastAPI với 20 biểu đồ tương tác cao.
      \item Tuân thủ chặt chẽ thiết kế UI/UX: đồng bộ bộ lọc trên URL, chế độ phóng to nhất quán.
    \end{itemize}
    \item \textbf{Tích hợp AI Human-in-the-loop:}
    \begin{itemize}
      \item Phát triển Trợ lý AI hỗ trợ phân tích tại chỗ bằng ngôn ngữ tự nhiên.
      \item Cơ chế ``chờ duyệt'' và thực thi cục bộ đảm bảo con người kiểm soát hoàn toàn hệ thống.
    \end{itemize}
  \end{itemize}
\end{frame}

% Slide 22
\begin{frame}[t]{Hạn chế}
  \begin{itemize}
    \setlength{\itemsep}{10pt}
    \item \textbf{Giới hạn về dữ liệu:}
    \begin{itemize}
      \item Khoảng trống số liệu (D7, D8 chỉ có từ 2018; một số địa phương thiếu dữ liệu).
      \item Giữ nguyên dữ liệu rỗng để bảo toàn độ tin cậy, làm giảm cỡ mẫu ở một số phép tính.
    \end{itemize}
    \item \textbf{Giới hạn phương pháp luận:}
    \begin{itemize}
      \item Phân tích định lượng (Pearson, K-Means) dừng ở mức mô tả độ phân tán và cùng biến thiên.
      \item Không đủ cơ sở để khẳng định mối quan hệ nhân quả.
    \end{itemize}
    \item \textbf{Giới hạn hệ thống AI:}
    \begin{itemize}
      \item Trình thực thi (Executor) đóng vai trò lớp bảo vệ cục bộ, chưa phải là sandbox công khai.
      \item Tính sẵn sàng phụ thuộc vào đường truyền mạng và quota của nhà cung cấp LLM (Groq).
    \end{itemize}
  \end{itemize}
\end{frame}

% Slide 23
\begin{frame}[t]{Hướng phát triển trong tương lai}
  \begin{itemize}
    \setlength{\itemsep}{12pt}
    \item \textbf{Mở rộng chiều phân tích dữ liệu:}
    \begin{itemize}
      \item Tích hợp chỉ số kinh tế\textendash xã hội địa phương (GRDP, tỷ lệ hộ nghèo, thu nhập bình quân).
      \item Xây dựng mô hình đa biến nhằm làm rõ tác động vĩ mô đến điểm quản trị công.
    \end{itemize}
    \item \textbf{Nâng cấp tính năng và bảo mật AI:}
    \begin{itemize}
      \item Cải thiện khả năng ghi nhớ ngữ cảnh dài hạn qua nhiều phiên thảo luận.
      \item Đưa Trình thực thi vào môi trường cô lập hoàn toàn (vd: Docker) nếu triển khai public.
    \end{itemize}
    \item \textbf{Đóng gói phần mềm ứng dụng:}
    \begin{itemize}
      \item Đóng gói Dashboard thành ứng dụng Desktop độc lập (standalone application).
      \item Hỗ trợ người dùng trải nghiệm ngay mà không cần tự cài đặt môi trường lập trình.
    \end{itemize}
  \end{itemize}
\end{frame}

% Slide 24
\begin{frame}[t]{Lời cảm ơn}
  \vspace{1cm}
  \begin{center}
    Nhóm xin gửi lời cảm ơn chân thành đến:\\
    \vspace{0.5cm}
    \textbf{Thầy Bùi Tiến Lên}\\
    \textbf{Thầy Trần Huy Bân}\\
    \textbf{Thầy Võ Nhật Tân}\\
    \vspace{0.5cm}
    đã tận tình giảng dạy, hướng dẫn và tạo điều kiện tốt nhất để nhóm hoàn thành đồ án môn học Trực quan hóa Dữ liệu.
    
    \vspace{1cm}
    \Large \textbf{Xin chân thành cảm ơn quý Thầy và các bạn đã chú ý lắng nghe!}
  \end{center}
\end{frame}
